\begin{thebibliography}{99}
\bibitem[\protect\citeauthoryear{Lipsitch et al.}{Lipsitch et al.}{2010}]{lipsitch2010}
Lipsitch, M., Tchetgen Tchetgen, E. J., and Cohen, T. 2010. Negative controls: A tool for detecting confounding and bias in observational studies. \textit{Epidemiology} 21(3), 383--388. \url{https://doi.org/10.1097/EDE.0b013e3181d61eeb}.

\bibitem[\protect\citeauthoryear{Fazil et al.}{Fazil et al.}{2024}]{fazil2024}
Fazil, M., Risquez, A., \& Halpin, C. (2024). A novel deep learning model for student performance prediction using engagement data. Journal of Learning Analytics, 11(2), 23-41. \url{https://doi.org/10.18608/jla.2024.7985}.

\bibitem[\protect\citeauthoryear{Ramaswami et al.}{Ramaswami et al.}{2023}]{ramaswami2023}
Ramaswami, G., Susnjak, T., \& Mathrani, A. (2023). Effectiveness of a learning analytics dashboard for increasing student engagement levels. Journal of Learning Analytics, 10(3), 115-134. \url{https://doi.org/10.18608/jla.2023.7935}.

\bibitem[\protect\citeauthoryear{Idowu et al.}{Idowu et al.}{2024}]{idowu2024}
Idowu, A., Koshiyama, A., \& Treleaven, P. (2024). Investigating algorithmic bias in student progress monitoring. Computers and Education: Artificial Intelligence, 7, 100267. \url{https://doi.org/10.1016/j.caeai.2024.100267}.

\bibitem[\protect\citeauthoryear{Palanci et al.}{Palanci et al.}{2024}]{palanci2024}
Palanci, A., Yilmaz, R. M., \& Turan, Z. (2024). Learning analytics in distance education: A systematic review study. Education and Information Technologies, 29, 22629-22650. \url{https://doi.org/10.1007/s10639-024-12737-5}.

\bibitem[\protect\citeauthoryear{Abdelrahman et al.}{Abdelrahman et al.}{2023}]{abdelrahman2023}
Abdelrahman, G., Wang, Q., \& Nunes, B. (2023). Knowledge tracing: A survey. ACM Computing Surveys, 55(11), 1-37. \url{https://doi.org/10.1145/3569576}.

\bibitem[\protect\citeauthoryear{Yang et al.}{Yang et al.}{2024}]{yang2024}
Yang, Q., Yang, Z., Cai, X., Zhao, H., Jia, J., \& Sun, F. (2024). Advances in methodologies of negative controls: A scoping review. Journal of Clinical Epidemiology, 166, 111228. \url{https://doi.org/10.1016/j.jclinepi.2023.111228}.

\bibitem[\protect\citeauthoryear{Tchetgen Tchetgen et al.}{Tchetgen Tchetgen et al.}{2024}]{proximal2024}
Tchetgen Tchetgen, E. J., Ying, A., Cui, Y., Shi, X., and Miao, W. 2024. An introduction to proximal causal inference. \textit{Statistical Science} 39(3), 375--390. \url{https://doi.org/10.1214/23-STS911}.

\bibitem[\protect\citeauthoryear{Zhang et al.}{Zhang et al.}{2025}]{zhang2025}
Zhang, D., Zhang, B., Wang, H., Lu, Y., Wolock, C. J., Hu, W., Wang, L., Hripcsak, G., \& Chen, Y. (2025). Negative control-calibrated difference-in-difference analyses: Addressing unmeasured confounding in real-world data with application to racial/ethnic differences. npj Digital Medicine, 8(1), 452. \url{https://doi.org/10.1038/s41746-025-01821-w}.

\bibitem[\protect\citeauthoryear{Tchetgen Tchetgen}{Tchetgen Tchetgen}{2014}]{coca2014}
Tchetgen Tchetgen, E. J. 2014. The control outcome calibration approach for causal inference with unobserved confounding. \textit{American Journal of Epidemiology} 179(5), 633--640. \url{https://doi.org/10.1093/aje/kwt303}.

\bibitem[\protect\citeauthoryear{Miao et al.}{Miao et al.}{2024}]{miao2024}
Miao, W., Shi, X., Li, Y., and Tchetgen Tchetgen, E. J. 2024. A confounding bridge approach for double negative control inference on causal effects. \textit{Statistical Theory and Related Fields} 8(4), 262--273. \url{https://doi.org/10.1080/24754269.2024.2390748}.

\bibitem[\protect\citeauthoryear{Park et al.}{Park et al.}{2024}]{park2024}
Park, C., Richardson, D. B., and Tchetgen Tchetgen, E. J. 2024. Single proxy control. \textit{Biometrics} 80(2), ujae027. \url{https://doi.org/10.1093/biomtc/ujae027}.

\bibitem[\protect\citeauthoryear{Schuemie et al.}{Schuemie et al.}{2018}]{schuemie2018}
Schuemie, M. J., Hripcsak, G., Ryan, P. B., Madigan, D., and Suchard, M. A. 2018. Empirical confidence interval calibration for population-level effect estimation studies in observational healthcare data. \textit{PNAS} 115(11), 2571--2577. \url{https://doi.org/10.1073/pnas.1708282114}.

\bibitem[\protect\citeauthoryear{Hwang et al.}{Hwang et al.}{2022}]{hwang2022}
Hwang, H., Quiroz, J. C., and Gallego, B. 2022. Assessing the effectiveness of empirical calibration under different bias scenarios. \textit{BMC Medical Research Methodology} 22, 208. \url{https://doi.org/10.1186/s12874-022-01687-6}.

\bibitem[\protect\citeauthoryear{Cobzaru et al.}{Cobzaru et al.}{2024}]{cobzaru2024}
Cobzaru, R., Welsch, R., Finkelstein, S., Ng, K., and Shahn, Z. 2024. Bias formulas for violations of proximal identification assumptions in a linear structural equation model. \textit{Journal of Causal Inference} 12(1), 20230039. \url{https://doi.org/10.1515/jci-2023-0039}.

\bibitem[\protect\citeauthoryear{Zheng et al.}{Zheng et al.}{2024}]{zheng2024}
Zheng, J., Wu, J., D'Amour, A., and Franks, A. 2024. Sensitivity to unobserved confounding in studies with factor-structured outcomes. \textit{Journal of the American Statistical Association} 119(547), 2026--2037. \url{https://doi.org/10.1080/01621459.2023.2240053}.

\bibitem[\protect\citeauthoryear{Kuzilek et al.}{Kuzilek et al.}{2017}]{kuzilek2017}
Kuzilek, J., Hlosta, M., and Zdrahal, Z. 2017. Open University Learning Analytics dataset. \textit{Scientific Data} 4, 170171. \url{https://doi.org/10.1038/sdata.2017.171}.

\bibitem[\protect\citeauthoryear{Weidlich et al.}{Weidlich et al.}{2022}]{weidlich2022}
Weidlich, J., Gasevic, D., \& Drachsler, H. (2022). Causal inference and bias in learning analytics: A primer on pitfalls using directed acyclic graphs. Journal of Learning Analytics, 9(3), 183-199. \url{https://doi.org/10.18608/jla.2022.7577}.

\bibitem[\protect\citeauthoryear{Cohausz}{Cohausz}{2022}]{cohausz2022}
Cohausz, L. (2022). When probabilities are not enough - A framework for causal explanations of student success models. Journal of Educational Data Mining, 14(3), 52-75. \url{https://doi.org/10.5281/zenodo.7304800}.

\bibitem[\protect\citeauthoryear{Inoue et al.}{Inoue et al.}{2025}]{inoue2025}
Inoue, K., Sakamaki, K., Komukai, S., Ito, Y., Goto, A., \& Shinozaki, T. (2025). Methodological tutorial series for epidemiological studies: Confounder selection and sensitivity analyses to unmeasured confounding from epidemiological and statistical perspectives. Journal of Epidemiology, 35(1), 3-10. \url{https://doi.org/10.2188/jea.JE20240082}.

\bibitem[\protect\citeauthoryear{Kennedy}{Kennedy}{2023}]{kennedy2023}
Kennedy, E. H. (2023). Towards optimal doubly robust estimation of heterogeneous causal effects. Electronic Journal of Statistics, 17(2), 3008-3049. \url{https://doi.org/10.1214/23-EJS2157}.

\bibitem[\protect\citeauthoryear{Inoue et al.}{Inoue et al.}{2024}]{inoue2024}
Inoue, K., Adomi, M., Efthimiou, O., Komura, T., Omae, K., Onishi, A., Tsutsumi, Y., Fujii, T., Kondo, N., \& Furukawa, T. A. (2024). Machine learning approaches to evaluate heterogeneous treatment effects in randomized controlled trials: A scoping review. Journal of Clinical Epidemiology, 176, 111538. \url{https://doi.org/10.1016/j.jclinepi.2024.111538}.

\bibitem[\protect\citeauthoryear{VanderWeele and Ding}{VanderWeele and Ding}{2017}]{vanderweele2017}
VanderWeele, T. J., \& Ding, P. (2017). Sensitivity analysis in observational research: Introducing the E-value. Annals of Internal Medicine, 167(4), 268-274. \url{https://doi.org/10.7326/M16-2607}.
\end{thebibliography}
